\section{Extended Interpretation}
The complete interpretation from the reference manuscript is retained below.
The main systems finding is that subscriber capacity and client access to that
capacity are separate. Adding Pods supplies service capacity; moving existing
subscription-bearing connections makes it reachable by persistent clients.
The guarded-greedy and DQL policies both achieve this in the scale-out model.
Their nearby throughput results support the shared mechanism rather than
requiring a claim that learning beats a strong greedy rule in every workload.

Policy choice depends on the operational objective. Greedy control gives a
direct, inspectable response to measured pressure. DQL can account for delayed
costs and several competing objectives. A subscription move changes future
backlog distribution; migration trades immediate handoff cost for later
capacity. Both selectors operate above the same bounded, checked actions.

QEdgeProxy selects service instances using client QoS feedback
\cite{Cilic2025}, while Yahya et al. adapt offloading through decentralized
bandits and server-reward adjustment \cite{Yahya2024}. The present work
focuses on subscription-aware persistent messaging clients under partial
observability and treats connection-level rebalancing as a first-class
control action alongside endpoint selection. This distinction
matters operationally: endpoint selection chooses the serving Pod, while
connection-level rebalancing corrects subscription skew within the client.

The holdout results show how learning compares with the heuristic.
Both policies can perform the same actions, so they can reach the same
capacity limit in some scenarios.  DQL's measured advantage is concentrated
in subscription distribution, combined connection imbalance, and avoiding a
harmful heuristic intervention during client slowdown.  The heuristic retains
an advantage in throughput and backlog-distribution metrics. These results
identify the operating objectives favored by each policy. The
component-level analysis reports these tradeoffs alongside the scalar reward;
policy selection uses development validation before the holdout is evaluated.

The distinction between a model's quality and the utility of the complete
system also appears in Casimiro et al.'s probabilistic framework for adapting
ML components \cite{Casimiro2024ModelChecking}. In our setting, a favorable
Q-value is therefore interpreted alongside backlog, throughput, and handoff
outcomes. Gheibi and Weyns introduce a lifelong-learning
layer to respond to changing adaptation options
\cite{Gheibi2024AdaptationDrift}. Our experiments evaluate a fixed policy under
the stated workload and deployment conditions. Lifelong adaptation is outside
this study's scope.

The loaded Oracle Cloud scale-out trial adds direct evidence for selective
execution (Section~\mainref{sec:live-dql-scaleout}). Every handoff preserved three
unselected active sessions, while the full cohort reached 8,000 ACK receipts.
The measured drain waits also show why opening a replacement socket is only
one part of migration: client processing and acknowledgement completion can
dominate the handoff duration. These measurements validate the learned
decision-to-execution path under the declared load and processing limit.

The action-contract studies provide a further design lesson: a useful
executor and a policy that reliably selects it are separate requirements.
The cap-restoring scale-up preflight improved reward and maximum subscription
count, but the expanded learner study did not pass every coverage gate.
Appendix~\ref{app:action-diagnostics} retains these results. They explain why
the live checkpoint uses a compatible, restricted rebalancing profile.

\subsection{Conclusion}

We presented a probe-guided, client-side framework for rebalancing
persistent subscription connections without requiring cluster-wide client
visibility.  The mechanism combines per-subscription queues, active-connection
telemetry, and a load-balancer-assigned probe. Greedy rules or DQL select
bounded actions, and a shared executor carries out receipt-gated handoff.
The technical contribution is the link between partial observations,
endpoint-preserving probe promotion, and selective subscription reassociation.

The primary experiment demonstrates recovery after a one-to-four pod capacity
increase with 20,000 subscriptions and 775 million queued events. Guarded-greedy
migration and frozen DQL deliver 221,260 and 219,762 events/s, respectively,
compared with 58,000 events/s for unchanged placement. Frozen DQL
uses all four pods, delivers 3.79 times the mean throughput, and leaves 27.52\%
less backlog after 20 minutes than retaining the ten original connection
assignments. Its 12.30\% reduction in backlog area satisfies the registered
paired statistical criterion. The adverse-probe ablation complements this
result: evaluating the candidate before promotion reduces backlog area by
4.19\% relative to forcing one harmful migration.

These findings identify the operational benefit of the mechanism: scaled-out
subscriber capacity becomes useful to existing subscriptions through
client-directed connection reassociation, while combined subscription and
backlog pressure determines whether a sampled target is suitable. The
supporting studies characterize policy tradeoffs and action coverage.
The mixed-workload holdout favors DQL in composite reward and subscription
balance, while the heuristic favors throughput and backlog balance. These
results support two evaluated policy options rather than universal DQL
superiority or an untested DQL-tuned scripted policy.
The isolated Oracle Cloud cluster qualification also demonstrates a 60\% reduction in source
subscriptions through three full-connection migrations, with all 2,000
subscriptions retained and 73,639 observed deliveries reconciled at closure.
Its startup measurements explicitly capture activation while deliveries flow.
A further loaded one-to-four-Pod trial retained all 1,000 subscriptions and
four active connections through three DQL-selected handoffs. Each action
preserved the other three sessions. The final client distribution was 500,
250, 250 and zero across the four available Pods, with 8,000 observed
deliveries acknowledged and no unresolved client deliveries at closure.
A separate Oracle Cloud preload pair measured twelve native TTL interruptions
and 89.44\% ACK-receipt completion, versus no interruptions and 100\% completion
in the TTL-disabled DQL arm (Section~\ref{sec:ttl-preload-results}). That arm
remained at guard-held \emph{stay}: a 200-second minimum telemetry sweep exceeded
the 180-second freshness deadline. The result measures expiry-related receipt
completion under the telemetry constraint. The comparative throughput and
backlog-recovery gains are simulator measurements under the stated service
assumptions. The Oracle Cloud experiments establish live handoff execution,
subscription retention, and observed delivery behavior. Together, these studies
describe the mechanism's measured benefits and the operating conditions under
which they were obtained.
